\begin{helper}
Let \(W\) be a finite vertex set, let \(v:\{1,\ldots,k\}\to W\) be an injective tuple, and let \(r:W\to\mathbb N\) be a ranking function that is injective on the range of \(v\). Among the \(k!\) permutations of the tuple positions, the permutations \(\sigma\) for which
\[
  r(v_{\sigma(1)})<r(v_{\sigma(2)})<\cdots<r(v_{\sigma(k)})
\]
form a one-element fiber. Equivalently,
\[
  k!\cdot
  \#\{\sigma\in\operatorname{Sym}(\{1,\ldots,k\}): r(v_{\sigma(1)})<\cdots<r(v_{\sigma(k)})\}
  =
  \#\operatorname{Sym}(\{1,\ldots,k\}).
\]
\end{helper}
\begin{proof}
Write \(f(i)=r(v_i)\). Since \(v\) is injective and \(r\) is injective on the range of \(v\), the values \(f(1),\ldots,f(k)\) are pairwise distinct.

First we prove uniqueness. Suppose \(\sigma\) and \(\tau\) are two permutations of the positions such that both sequences
\[
  f(\sigma(1)),\ldots,f(\sigma(k))
  \quad\text{and}\quad
  f(\tau(1)),\ldots,f(\tau(k))
\]
are strictly increasing. Each sequence lists the same set \(\{f(1),\ldots,f(k)\}\), because \(\sigma\) and \(\tau\) are permutations. A finite set of distinct integers has a unique increasing listing: the first term is its least element, and after removing the first \(m\) least elements, the next term is the least remaining element. Hence \(f(\sigma(i))=f(\tau(i))\) for every \(i\). Since the \(f(i)\) are pairwise distinct, \(\sigma(i)=\tau(i)\) for every \(i\), so \(\sigma=\tau\). Thus there is at most one increasing permutation.

Now list the set \(\{f(1),\ldots,f(k)\}\) in increasing order as
\[
  \rho_1<\rho_2<\cdots<\rho_k .
\]
For each \(i\), choose the unique position \(\sigma(i)\) with \(f(\sigma(i))=\rho_i\). This defines a permutation of the \(k\) positions, because every position has one of the ranks \(\rho_i\) and the ranks are distinct. By construction \(f(\sigma(1))<\cdots<f(\sigma(k))\), so an increasing permutation exists.

Therefore the set of increasing permutations has cardinality exactly \(1\). The total number of permutations of \(k\) positions is \(k!\). Multiplying the one-element fiber by \(k!\) gives exactly the total number of permutations, which is the claimed identity.
\end{proof}
